%=====================================================================
% ICRA-style METHOD section for the RH-NBV planner (simplified).
% For the IEEEtran conference template (two-column):
%   \documentclass[conference]{IEEEtran}
%   \usepackage{amsmath,algorithm,algpseudocode,booktabs,cite}
% Citation keys match references.bib in this repo.
%=====================================================================

\section{PROBLEM FORMULATION}
\label{sec:problem}

A 6-DoF manipulator carries an RGB-D camera and must perceive a target
object at a known approximate position $\mathbf{t}$, despite occlusion.
The scene belief is a voxel grid $M$ that stores, for every voxel, an
occupancy probability $p_o$ and a semantic probability $p_s$
\cite{burusa2024}. The camera moves inside a workspace $\mathcal{B}$ and
always looks at $\mathbf{t}$, so a viewpoint is simply a position
$\xi \in \mathcal{B}$. Given a budget of $N$ views, the goal is to
reconstruct the region of interest (ROI) around $\mathbf{t}$ as well as
possible while moving the camera as little as possible.

\section{PROPOSED APPROACH}
\label{sec:method}

The planner repeats four steps (Algorithm~\ref{alg:rh}): sample $K$
candidate view sequences of length $H$, score each sequence, move to the
first view of the best sequence, and replan after the new measurement.
Planning $H$ steps ahead avoids the myopia of single-step NBV
\cite{burusa2024}; executing only the first step keeps the planner
robust, because every new measurement can change what the best plan is
\cite{bircher2016}.

\subsection{Viewpoint Utility}
\label{sec:method-utility}

How much is a single view worth? We cast rays from the candidate
position through the voxel grid. Each sampled point $i$ on a ray
contributes its semantic entropy $\mathcal{H}_s(i)$ (how uncertain the
grid still is about that point), weighted by the transmittance
$T(i) = \prod_{j<i}(1 - o(j))$, i.e.\ the probability that the ray
actually reaches the point without being blocked
\cite{delmerico2018}. Occluded voxels therefore contribute almost
nothing. The utility of a view is the log of the mean contribution:
%
\begin{equation}
  f(M, \xi) \;=\; \log \Bigl(
  \tfrac{1}{|\mathcal{R}| N_r}
  \textstyle\sum_{r \in \mathcal{R}} \sum_{i=1}^{N_r}
  T(i)\, \mathcal{H}_s(i) + \epsilon \Bigr).
  \label{eq:utility}
\end{equation}
%
The log does not change the ranking of candidates; it only brings the
gain to a scale on which one fixed motion-cost weight $\lambda$ works
for all iterations.

\subsection{Generating Candidate Sequences}
\label{sec:method-candidates}

A sequence is grown step by step from the current position. With
probability $b$ the next view is a small \emph{orbit step}: the camera
slides tangentially on the sphere around $\mathbf{t}$. Since the camera
always faces the target, only tangential motion changes the viewing
angle---and new information comes from new angles. With probability
$1-b$ the next view is a uniform random jump inside $\mathcal{B}$, which
lets the planner consider completely different regions. Every generated
position is pushed back to a minimum stand-off distance from the target
and clamped to the robot reach, so all candidates are feasible.

\subsection{Scoring a Sequence}
\label{sec:method-evaluation}

A sequence $\Xi = (\xi_0, \dots, \xi_{H-1})$ is scored by
%
\begin{equation}
  J(\Xi) \;=\; \sum_{k=0}^{H-1} \gamma^{k}
  \Bigl[ \rho_k\, f\bigl(M_{\mathrm{pred}}, \xi_k\bigr)
  \;-\; \lambda\, \lVert \xi_k - \xi_{k-1} \rVert_2 \Bigr].
  \label{eq:objective}
\end{equation}
%
Three ingredients matter here:

\emph{a) Forward simulation.} The gain of step $k$ is computed on a
predicted map $M_{\mathrm{pred}}$, not on the current map. After each
hypothetical view, the uncertain voxels it would see are moved partially
towards confident values. This way, step 2 gets no credit for
information that step 1 would already have collected. Without this, two
nearby views would be rewarded twice for the same unknown region, which
is exactly what happens in the additive gain of
\cite{bircher2016, zhang2022}.

\emph{b) Redundancy penalty.} If a view lies closer than $s/2$ to an
earlier view of the same sequence, its gain is multiplied by
$\rho_k = 0.3$ (otherwise $\rho_k = 1$). This discourages sequences that
barely move.

\emph{c) Discount and motion cost.} $\gamma < 1$ trusts later steps
less, since they rely on predictions of predictions, and $\lambda$
penalises long moves.

\subsection{Stagnation Escape}
\label{sec:method-stagnation}

Orbit steps are local, so the planner can get stuck once a region is
exhausted. If ROI coverage improves by less than $\tau$ for $P$
consecutive iterations, the camera jumps to the farthest of $20$ random
reachable positions. This handles the premature-convergence problem
reported for purely local planners in occluded scenes
\cite{verlinden2024pso}.

\begin{algorithm}[t]
\caption{RH--NBV: one planning iteration}
\label{alg:rh}
\begin{algorithmic}[1]
\Require current position $\xi_0$, belief $M$, target $\mathbf{t}$
\If{coverage stagnated for $P$ iterations}
    \State jump to a distant reachable position
\EndIf
\For{$i = 1, \dots, K$}
    \State $\Xi_i \gets$ grow $H$ steps (orbit step with prob.~$b$,
           else random jump); keep positions feasible
    \State $J_i \gets$ score $\Xi_i$ with Eq.~\eqref{eq:objective}
\EndFor
\State $\Xi^{*} \gets \arg\max_i J_i$
\State \Return first view of $\Xi^{*}$, looking at $\mathbf{t}$
\end{algorithmic}
\end{algorithm}

\subsection{Implementation Details}
\label{sec:method-implementation}

We use $H{=}3$, $K{=}10$, $\gamma{=}0.85$, $\lambda{=}2.0$, $b{=}0.7$,
$P{=}4$, $\tau{=}1.5\%$, and $s{=}0.065$\,m (the step size of
\cite{burusa2024}); $H$ is varied in an ablation over $\{1,2,3,4\}$.
One iteration performs $K \cdot H = 30$ gain evaluations on the GPU.
Following the cost metric of \cite{burusa2024}, the forward-simulation
updates are not counted as ray-tracing calls, but their runtime is
included in all timing results. For fairness, the two RH-specific
extensions with no counterpart in GradientNBV (a spherical orbit shell
and an occlusion-bonus term) are disabled in all baseline comparisons
and studied only as ablations; both planners use identical voxel grids,
workspaces, ROIs, and random seeds.
